% !TeX root = ../tesis.tex
%%%%%%%%%%%%%%%%%%%%%%%%%%%%%%%%%%%%%%%%%%%%%%%%%%%%%%%%%%%%%%%%%%%
% SECCIÓN 6.5 — CONCLUSIONES DEL ESCENARIO
% Issue #25: conclusiones del Cap. 6 y caso más favorable (tarea del
% asesor). Lectura urbana de §6.3–§6.4; solo citas ya presentes en la
% tesis. No usa el hallazgo especulativo de transbordos (CONTEXTO_CAP6).
%
% Estructura:
%   Entrada
%   §6.5.1  Acercar no es conectar: el anillo en los contornos
%   §6.5.2  De red radial a red anillo-radial
%           (+ Programa Maestro del Metro: Década de 1980 y
%              Horizonte 2010; ddf1985programa)
%   §6.5.3  Los nodos de transbordo como CETRAM de nueva generación
%   §6.5.4  Modo y derecho de vía: el caso más favorable
%   §6.5.5  Síntesis (tab:anillo-sintesis; remite a las figuras de T,
%           B(v) y banda dominante. El mapa de síntesis se descartó para la
%           versión Alfa, 2026-10-25)
% Corrida: 2026-10-04. T por demarcación sobre los 10,610 nodos comunes
%   (VFTModel/notebooks/ASSETS/t_distribucion/t_comparativo_origen.csv).
%%%%%%%%%%%%%%%%%%%%%%%%%%%%%%%%%%%%%%%%%%%%%%%%%%%%%%%%%%%%%%%%%%%
\section{Conclusiones del Escenario}
\label{sec:anillo-conclusiones}

Con el Espectro MJG de los tres casos termina la parte cuantitativa del
capítulo. Queda decir qué significan esos números para la ciudad, y a eso se
dedica esta sección. Primero se revisa a qué contornos de la metrópoli llega el
anillo; luego, cómo cambia la dependencia de la red respecto de sus nodos
centrales; después, qué función tendrían las estaciones de transbordo del
anillo, y al final, cuál de los dos casos resulta más favorable y por qué. Un
cuadro de síntesis cierra la sección. Conviene recordar que la lectura es
estructural: el modelo dice por dónde pasan los caminos más cortos y cuánto
duran, no cuántas personas los usarían ni cuánto costaría construir el corredor.

\subsection{Acercar no es Conectar: el Anillo en los Contornos}
\label{subsec:anillo-conclusiones-contornos}

La primera conclusión es que el anillo casi no acerca la red a territorio nuevo.
La cobertura pasa de 13.29 a 13.43\,\% de la \zmvm{} en los dos casos
(§\ref{subsec:anillo-cobertura}), porque el Periférico corre por la parte de la
ciudad que ya tiene más estaciones. Los 41 municipios que hoy no tienen ninguna
parada a distancia caminable siguen igual; están en la periferia más alejada
del centro, lejos del trazo (Figura~\ref{fig:semovi-contornos}). Nada cambia
para ellos. Un corredor sobre el Periférico atiende a quien
ya vive cerca de la red, no a quien vive lejos de ella.

Lo que el anillo sí cambia es la conexión de quienes ya tienen estaciones cerca.
Como se observa en la Figura~\ref{fig:anillo-t-demarcacion}, los municipios del
oriente metropolitano, Nezahualcóyotl, Chimalhuacán y Ecatepec, ganan entre
media hora y una hora en el tiempo medio de sus nodos, tanto con el Anillo MB
como con el Anillo Metro. En la red actual son de los más apartados del resto de
la ciudad, y el anillo les da un camino lateral que antes no tenían. Con el
Anillo Metro se suma el sur. Tlalpan, Xochimilco, Tláhuac y La Magdalena
Contreras, que con el Anillo MB ganan seis o siete minutos, pasan a ganar entre
15 y 27.

Este resultado responde en parte a la segunda pregunta secundaria, que
vinculaba la falta de conexión lateral en el Tercer Contorno con los tiempos de
traslado de quienes vienen del Cuarto (§\ref{subsec:pregunta}). En el modelo,
esa falta sí alarga los viajes de la periferia: cuando el anillo ofrece un
enlace lateral, los tiempos bajan más donde la red actual ofrecía menos
caminos. La mejora depende del modo. A la velocidad del Mexibús llega sobre todo
al oriente; solo a la del Metro alcanza también a las alcaldías del sur que el
Capítulo~1 identifica con el Tercer Contorno (§\ref{subsec:beneficiarios-contornos}).
Queda sin responder cuántas personas viven cerca de esos nodos, porque $T$ no
pondera por población.

\subsection{De Red Radial a Red Anillo-Radial}
\label{subsec:anillo-conclusiones-centralidad}

La primera pregunta secundaria se refería a la dependencia de la red respecto de
unos cuantos nodos centrales (§\ref{subsec:pregunta}). En la red actual esa
dependencia tiene nombre: Tacubaya concentra la mayor intermediación, seguida de
Mixcoac e Hidalgo, y los nodos con la banda dominante más alta dibujan un
corredor de Garibaldi a Barranca del Muerto (§\ref{subsec:anillo-espectro-nodo}).
El Anillo MB no cambia ese cuadro: Tacubaya pierde un 3.1\,\% y los diez
primeros nodos siguen en el centro. Con el Anillo Metro es otra cosa. La red
deja de tener un solo eje, porque Tacubaya pierde casi un tercio de su
intermediación y nueve de los diez primeros nodos pasan al Periférico, en
Iztapalapa, Álvaro Obregón, Xochimilco y La Magdalena Contreras
(Figura~\ref{fig:anillo-bv-alcaldias}).

Que la intermediación salga del centro no significa que se reparta. El Gini sube
de 0.8007 a 0.8134 con el Anillo Metro, ya que los caminos que antes pasaban por
Tacubaya o Mixcoac pasan ahora por unas cuantas estaciones del Periférico. La
red deja de ser solo radial. Se vuelve anillo-radial, con un eje central y otro
sobre el arco sur y oriente del anillo, y tener dos ejes en lugar de uno podría
hacerla menos frágil ante una falla en el centro. Pero el modelo no simula
cierres ni contingencias; esa ventaja es por ahora una hipótesis.

Los dos referentes revisados en la sección~\ref{sec:casos-similares} ayudan a
leer esta diferencia. La Gran Línea Circular de Moscú se pensó para descargar
las líneas radiales del centro \citep{efe2023}, y es lo que el Anillo Metro hace
en el modelo con Tacubaya, Mixcoac y la Línea~2. El \termino{Suburban Rail Loop}
de Melbourne busca unir entre sí las zonas del \termino{Middle-Ring} sin pasar
por el centro \citep{suburbanrailloop2021}, y en la \zmvm{} son el sur y el
oriente los que ganan intermediación. Los dos proyectos son ferroviarios. Que
aquí solo el caso con la velocidad del Metro reproduzca su efecto es coherente
con esa elección, aunque se trata de un resultado del modelo y no de una
evaluación de esos sistemas.

% Programa Maestro del Metro, segunda revisión, versión 1985
% (DDF / COVITUR). Fotografías en Figures/Cap6/Planes-Maestros-Previos/.
% Las fotos de PlanMaestro2010/ son el horizonte 2010 del mismo documento
% (pp. 63--66), no un plan de la década de 2010.
\subsubsection{El Programa Maestro del Metro: Década de 1980}
\label{subsubsec:anillo-pm-1980}

La idea de un transporte masivo sobre el Periférico tiene un antecedente en la
propia ciudad. El Programa Maestro del Metro, elaborado en la década de 1980 por
el \textsc{ddf} a través de la Comisión de Vialidad y Transporte Urbano
(\textsc{covitur}), con una segunda revisión de 1985 \citep{ddf1985programa},
dibujó redes de 444~km en 1982 y de 416~km en 1983. En las dos aparecen líneas
sobre el Periférico en el sur y el oriente
(Figura~\ref{fig:anillo-pm-1982-1983}). El \textsc{api-i} retoma, así, un trazo
que la planeación del Metro ya había considerado sobre la vialidad que dejó la
regencia de Uruchurtu (§\ref{sec:anillo-periferico}). Aquí el programa sirve
como antecedente y no como fuente de datos del modelo.

\begin{figure}[p]
    \centering
    \includegraphics[angle=90,height=0.82\textheight]{%
        Figures/Cap6/Planes-Maestros-Previos/PlanMaestro1983-1986/PlanMaestro1}
    \caption[Planes Maestros del Metro de 1982 y 1983]{%
    Redes del Metro de los Planes Maestros de 1982 (444~km, esquema E-5) y de
    1983 (416~km, esquema E-6). Fotografía del documento original. Otra toma
    en la Figura~\ref{fig:anx-g-pm-1982-1983}.}
    \label{fig:anillo-pm-1982-1983}
    \fuentefigura{\citet[p.~23]{ddf1985programa}; fotografía del autor.}
\end{figure}

Para estimar la demanda, el programa usó la encuesta Origen-Destino de 1983 y un
modelo en cuatro pasos (generación, atracción, distribución y asignación de
viajes), con la ciudad dividida en tres regiones
\citep[pp.~54--58]{ddf1985programa}. La Región Centro quedaba dentro del
Circuito Interior; la Corona, entre el Circuito Interior y el Periférico, y la
Exterior, del Periférico hacia afuera. Es decir, en 1983 el Periférico marcaba
el paso de la ciudad central a su periferia. De la asignación salieron 34
corredores de movimiento, que después se jerarquizaron con un modelo de
evaluación. En ese orden, el corredor de la Vía Tapo al Periférico Sur quedó en
el lugar 17 y el de la Calzada de las Armas por el Periférico Norte, Centro y
Sur, en el 33 (Cuadro~\ref{tab:anillo-pm-corredores}).

\begin{table}[H]
    \centering
    \caption[Corredores de movimiento del Programa Maestro del Metro]{%
    Jerarquización de los 34 corredores de movimiento del Programa Maestro del
    Metro (Tabla T-5 del original). En negritas, los dos corredores sobre el
    Periférico. Transcripción con ortografía actualizada.}
    \label{tab:anillo-pm-corredores}
    \small
    \setlength{\tabcolsep}{4pt}
    \begin{tabularx}{\textwidth}{|r|X|r|X|}
        \hline
        \textbf{Lugar} & \textbf{Corredor} & \textbf{Lugar} & \textbf{Corredor} \\
        \hline
        1  & Av. Chapultepec -- Ignacio Zaragoza &
        18 & Río de la Magdalena -- Sta. Ana -- Calz. Tulyehualco \\
        2  & México-Tacuba -- Tlalpan &
        19 & Laguna de Términos -- Rep. de Uruguay \\
        3  & Insurgentes Nte. -- Cuauhtémoc -- Universidad &
        20 & Av. Gran Canal -- Río Churubusco -- Cafetales \\
        4  & Inguarán -- Fco. Morazán -- La Viga &
        21 & Miguel Á. de Quevedo -- Calz. Tasqueña \\
        5  & Lázaro Cárdenas -- Río Consulado -- Hangares &
        22 & Félix Cuevas -- Municipio Libre -- A. de Circunvalación \\
        6  & Calz. Guadalupe -- Calz. Azcapotzalco-La Villa -- San Juan de Aragón &
        23 & Gral. Emiliano Zapata -- Ote. 172 \\
        7  & Aquiles Serdán -- Av. Revolución &
        24 & Ángel Urraza -- Trabajadoras Sociales \\
        8  & Calz. Guadalupe -- Brasil -- Amp. Plutarco E. Calles &
        25 & Benjamín Franklin -- Plutarco E. Calles \\
        9  & Baja California -- Chabacano -- Av. Río de la Piedad &
        26 & Dr. Río de la Loza -- Fray Servando T. de Mier -- Av. 8 \\
        10 & Blvd. Miguel C. Saavedra -- Mosqueta -- Oceanía &
        27 & Av. Cuitláhuac -- Alfredo Robles D. \\
        11 & Calz. Vallejo -- Lázaro Cárdenas -- Div. del Norte -- San Benjamín &
        28 & Deportivo Reynosa -- Montevideo \\
        12 & Popocatépetl -- Ermita Iztapalapa &
        29 & Calle 4 -- Ote. 253 -- Calle 10 -- San Fco. de Jesús \\
        13 & Eugenia -- M. Buendía &
        30 & Eduardo Molina -- Fco. del Paso -- Av. 5 -- Cafetales \\
        14 & Rojo Gómez -- Calz. San Lorenzo &
        31 & Centenario -- Calz. de la Viga -- Circunvalación -- Miramontes \\
        15 & Heliópolis -- Manuel González -- Transvaal &
        32 & Mariano Escobedo -- Div. del Norte \\
        16 & Av. Ceylán -- Insurgentes Norte, Centro, Sur &
        \textbf{33} & \textbf{Calz. Las Armas -- Periférico Norte, Centro, Sur} \\
        \textbf{17} & \textbf{Vía Tapo (Proy. Anillo Periférico) -- Periférico Sur} &
        34 & Ignacio Zaragoza \\
        \hline
    \end{tabularx}
    \fuentefigura{Elaboración propia con base en \citet{ddf1985programa},
    Tabla T-5.}
\end{table}

El modelo de evaluación combinó nueve parámetros, cada uno con un peso fijado
por consenso de un grupo de personas de distintas actividades y estratos
socioeconómicos (Cuadro~\ref{tab:anillo-pm-parametros}). La captación de
pasajeros, global y por kilómetro de ampliación, suma casi la cuarta parte del
peso. Los demás parámetros también salen de la demanda de 1983 asignada a la
red: las horas-hombre de viaje, la congestión del Metro y de los autobuses, el
equilibrio de las cargas o la población beneficiada. Ninguno mide qué tanto
depende la red de unos cuantos nodos ni qué tan unidas quedan entre sí las
zonas de la periferia. El más cercano a esa idea, el aumento de la redundancia,
pesa 9.33\,\%.

\begin{table}[H]
    \centering
    \caption[Parámetros de evaluación del Programa Maestro del Metro]{%
    Parámetros de evaluación de los corredores del Programa Maestro del Metro
    y su peso relativo (Tabla T-4 del original). Transcripción con ortografía
    actualizada.}
    \label{tab:anillo-pm-parametros}
    \small
    \setlength{\tabcolsep}{4pt}
    \begin{tabularx}{\textwidth}{|p{4.6cm}|X|r|}
        \hline
        \textbf{Parámetro} & \textbf{Descripción} & \textbf{Peso (\%)} \\
        \hline
        Captación global del Metro & Suma de las captaciones de todas las
            líneas del sistema & 11.43 \\
        \hline
        Captación específica de la ampliación & Nuevos pasajeros por cada
            kilómetro de ampliación & 12.54 \\
        \hline
        Horas-hombre en el sistema & Tiempo empleado en transportarse por toda
            la población & 14.57 \\
        \hline
        Nivel de congestión del sistema de superficie & Nivel de ocupación de la
            red de autobuses & 11.53 \\
        \hline
        Congestión del sistema & Congestionamiento en las estaciones del Metro
            & 11.66 \\
        \hline
        Equilibrio de las cargas & Estado de la carga por tramos con relación a
            la capacidad de las líneas del sistema & 9.90 \\
        \hline
        Congruencia con el Programa de Desarrollo Urbano del D.F. &
            Conectabilidad de las líneas con los centros urbanos del Programa de
            Desarrollo Urbano del D.F. & 8.72 \\
        \hline
        Población beneficiada & Población beneficiada directamente por las
            ampliaciones & 10.32 \\
        \hline
        Incremento en el nivel de redundancia & Diversificación de los
            recorridos realizables por Metro & 9.33 \\
        \hline
        \multicolumn{2}{|r|}{\textbf{Total}} & \textbf{100.00} \\
        \hline
    \end{tabularx}
    \fuentefigura{Elaboración propia con base en
    \citet[p.~60]{ddf1985programa}, Tabla T-4.}
\end{table}

\subsubsection{El Programa Maestro del Metro: Horizonte 2010}
\label{subsubsec:anillo-pm-2010}

Con esos corredores, el programa definió la red del Metro para el horizonte
2010. El criterio fue el de una «red equilibrada», en la que la oferta de cada
línea correspondiera a su demanda en la hora de máxima demanda, para evitar
líneas saturadas o subutilizadas \citep[p.~63]{ddf1985programa}. Las 34 líneas
iniciales sumaban 655~km y ofrecían 6.23 millones de viajes en esa hora, frente
a una demanda estimada de 2.20 millones. Para reducir el exceso de oferta, la
red se simuló varias veces y en cada vuelta se quitaron los corredores que no se
justificaban por captación (Cuadro~\ref{tab:anillo-pm-depuracion}). En la
tercera vuelta salió el Periférico Poniente y en la cuarta, el Periférico
Oriente. La red final quedó en 15 líneas y 315~km.

\begin{table}[H]
    \centering
    \caption[Depuración de la red del Metro al horizonte 2010]{%
    Depuración de la red del Metro al horizonte 2010 en el Programa Maestro:
    líneas, longitud y corredores eliminados en cada simulación. En negritas,
    los corredores del Periférico.}
    \label{tab:anillo-pm-depuracion}
    \small
    \setlength{\tabcolsep}{4pt}
    \begin{tabularx}{\textwidth}{|l|r|r|X|}
        \hline
        \textbf{Paso} & \textbf{Líneas} & \textbf{km} &
            \textbf{Corredores eliminados o fusionados} \\
        \hline
        Red inicial (E-12) & 34 & 655 & --- \\
        \hline
        Simulación 1 & 26 & 563 & División del Norte; Ejes Viales 1 Oriente,
            5 Norte, 4 Sur, 6 Sur y 7-A Sur; fusión de las líneas 19, 26 y 34 \\
        \hline
        Simulación 2 & 20 & 454 & Eduardo Molina, Fco. del Paso y Troncoso,
            Oriente 253; Ejes Viales 5 Oriente, 2 Norte, 7 Sur y 10 Sur \\
        \hline
        Simulación 3 & 18 & 397 & \textbf{Periférico Poniente}; Eje Central
            Lázaro Cárdenas \\
        \hline
        Simulación 4 & 16 & 345 & \textbf{Periférico Oriente};
            Cuitláhuac--Robles Domínguez \\
        \hline
        Red 2010 (E-13) & 15 & 315 & Línea 19 y prolongaciones de las líneas
            4, 12 y 20 \\
        \hline
    \end{tabularx}
    \fuentefigura{Elaboración propia con base en
    \citet[p.~63]{ddf1985programa}.}
\end{table}

Hubo además un límite administrativo. La red se diseñó dentro del Distrito
Federal, con la excepción de las Líneas~2 y~14, aunque el propio documento
reconoce que la demanda habría justificado líneas en el Estado de México
\citep[p.~66]{ddf1985programa}. Para compensarlo, las líneas que llegaban al
límite tendrían terrenos para centros de intercambio modal que las unieran con
el transporte del resto del área metropolitana, una idea muy parecida a la de
§\ref{subsec:anillo-conclusiones-cetram}. Solo cambia el lugar. En 1985 el
intercambio se pensaba en la frontera política y no sobre el anillo.

Leído desde esta tesis, el programa ayuda a entender por qué el anillo no se
construyó. Con la demanda que registró la encuesta de 1983, los corredores del
Periférico no alcanzaban la captación que exigía una red equilibrada, y la
evaluación no tenía forma de medir lo que el anillo aportaba a la estructura de
la red. En las décadas siguientes la \zmvm{} creció más allá del Periférico de
forma discontinua, con el patrón saltón de la Figura~\ref{fig:sprawl}, y lo
hizo sin un enlace lateral entre sus nuevas zonas. Hoy la red que diagnosticó
el Capítulo~\ref{ch:analisis-red-actual} depende del centro para casi cualquier
viaje largo. El corredor anillar no es, por eso, una ocurrencia. Es una idea
que la planeación del Metro ya tenía en los años ochenta y que se descartó con
los indicadores de su época.

Esto no significa que la evaluación de 1985 estuviera mal ni que esta tesis la
sustituya. Aquella midió con detalle la demanda y no la estructura de la red;
esta mide la estructura y deja fuera la demanda
(§\ref{subsec:corredor-anillar-limitaciones}). Lo que aporta la comparación es
que, a la velocidad del Metro, el anillo acorta los viajes de la periferia y
quita peso al centro, dos efectos que la captación por corredor no podía
registrar. La pregunta que llevó a descartarlo, cuánta gente lo usaría, sigue
abierta y requiere juntar ambos enfoques.

\subsection{Los Nodos de Transbordo como \textsc{cetram} de Nueva Generación}
\label{subsec:anillo-conclusiones-cetram}

Con el Anillo Metro aparece un tipo de nodo que la red actual no tiene.
Estaciones como Periférico Oriente/Tláhuac, Constitución de 1917 y Canal
Nacional, en el oriente y el sur, o Rómulo O'Farril, Barranca del Muerto y Luis
Cabrera, en el poniente, quedan con una intermediación de 0.12 a 0.16, cercana a
la de Tacubaya. Su Fuerza Capilar, en cambio, va de 6 a 12 conexiones,
alrededor de la mediana de la red, que es de 7. Son nodos por los que pasarían
muchos caminos, pero a los que llegan pocas rutas de superficie. En la red
actual pasa lo contrario con los \textsc{cetram} formales, que reúnen muchas
rutas y quedan lejos de los primeros lugares de intermediación
(§\ref{subsubsec:glosario-hub-cetram}).

Si el anillo se construyera, estas estaciones tendrían que planearse como
centros de transferencia y no como paradas de paso. Eso implica rutas
alimentadoras que las unan con las colonias cercanas, espacio para el transbordo
entre la \textsc{rtp}, los Corredores Concesionados y el anillo, y accesos de
última milla para quien llega a pie o en bicicleta. Serían \textsc{cetram} de
nueva generación. A diferencia de Pantitlán o Indios Verdes, que crecieron
alrededor de terminales de líneas radiales, estarían en la periferia y se
pensarían desde el inicio para el viaje lateral. Se trata de una recomendación
que se desprende de la estructura de la red; cuántas personas usarían cada
estación y qué suelo hay disponible en cada sitio son preguntas que el modelo no
responde (§\ref{subsec:corredor-anillar-limitaciones}).

\subsection{Modo y Derecho de Vía: el Caso más Favorable}
\label{subsec:anillo-conclusiones-modo}

Los Casos~2 y~3 comparten trazo, estaciones y penalización por transbordo; solo
cambian la velocidad y el derecho de vía (§\ref{subsec:anillo-parametros}). Esa
diferencia basta para separar sus resultados. El Anillo MB, a unos 14~km/h, le
gana a la \textsc{rtp} que hoy circula por el Periférico, pero no al Metro del
centro, así que toma el papel de la \textsc{rtp} sin cambiar la forma de la red.
El Anillo Metro, a 36~km/h, compite con las líneas radiales y por eso cambia la
jerarquía de la red: baja más el tiempo de viaje, alcanza al sur y al poniente y
genera los nodos de transbordo descritos arriba. Por estas razones, el Caso~3 es
el más favorable de los dos.

La conclusión tiene también una lectura histórica. El Periférico se construyó
para el automóvil, dentro de un programa de obras viales que se impuso a otras
necesidades de la ciudad (§\ref{subsec:sesgo-vial}), y en esos mismos años el
Metro se rechazó por costoso e innecesario (§\ref{subsec:rechazo-metro}). Los
resultados del Caso~3 sugieren que la vialidad hecha para el coche es la que
mejor podría alojar un transporte masivo lateral, si cuenta con vía exclusiva y
una velocidad parecida a la del Metro. Con un \brt{} confinado, el anillo acorta
los viajes, pero repite lo que ya hace la red de superficie. La comparación no
incluye el costo de construir y operar cada opción.

\subsection{Síntesis}
\label{subsec:anillo-conclusiones-sintesis}

El Cuadro~\ref{tab:anillo-sintesis} resume la lectura de los dos casos frente a
la red actual en los cuatro puntos de esta sección. Su expresión en el
territorio ya se mostró antes: el cambio del tiempo de viaje por demarcación
(Figura~\ref{fig:anillo-t-demarcacion}), el reparto de la intermediación
(Figura~\ref{fig:anillo-bv-comparativo}) y los nodos que cambian de banda
dominante (Figura~\ref{fig:anillo-banda-cambio}).

\begin{table}[H]
    \centering
    \caption[Síntesis de los Casos~2 y~3]{%
    Síntesis de la lectura urbana de los Casos~2 y~3 frente a la red actual.
    El cambio de $T$ corresponde a los 10,610 nodos comunes.
    Corrida: 2026-10-04.}
    \label{tab:anillo-sintesis}
    \small
    \setlength{\tabcolsep}{4pt}
    \begin{tabularx}{\textwidth}{|l|X|X|}
        \hline
        \textbf{Aspecto} & \textbf{Caso~2 (Anillo MB)} & \textbf{Caso~3 (Anillo Metro)} \\
        \hline
        Cobertura ($C$) & 13.43\,\% de la \zmvm{} (13.29 en el Caso~1); casi sin
            territorio nuevo & Igual que el Caso~2 \\
        \hline
        Tiempo de viaje ($T$) & $-4.1$\,\%; la mejora se concentra en el oriente
            metropolitano & $-10.1$\,\%; mejoran el oriente, el sur y el poniente \\
        \hline
        Intermediación ($B(v)$) & Tacubaya $-3.1$\,\%; los diez primeros nodos
            siguen en el centro; Gini 0.8012 & Tacubaya $-29.7$\,\%; nueve de los
            diez primeros nodos están en el anillo; Gini 0.8134 \\
        \hline
        Forma de la red & Radial; el anillo releva a la \textsc{rtp} del
            Periférico & Anillo-radial, con un eje central y otro en el arco sur
            y oriente \\
        \hline
        Nodos de transbordo & Sin cambios relevantes & Seis estaciones con $B(v)$
            de 0.12 a 0.16 y de 6 a 12 conexiones \\
        \hline
        Espectro MJG & Solo mejora $T$ (0.748); 10 nodos cambian de banda
            dominante & $T$ pasa a 0.75--1.00 (0.822); 21 nodos cambian de banda
            dominante \\
        \hline
        Lectura & Mejora local que conserva la estructura de la red & Caso más
            favorable: cambia la estructura, aunque no reduce la concentración \\
        \hline
    \end{tabularx}
    \fuentefigura{Elaboración propia con base en §\ref{sec:anillo-indicadores}
    y §\ref{sec:anillo-espectro}.}
\end{table}

En conjunto, el anillo cumple solo una parte de lo que se esperaba de él. Acorta
los viajes en toda la red y, a la velocidad del Metro, saca del centro una parte
de la intermediación. No amplía la cobertura. Tampoco reparte la dependencia
entre más nodos, porque un solo corredor sobre el Periférico no llega a
territorio nuevo y junta en sus estaciones los caminos que deja el centro. El
Capítulo~\ref{ch:conclusiones} retoma estos resultados para contrastarlos con la
pregunta de investigación y la hipótesis.
